\section*{Acknowledgements and Declaration of Authorship}
\addcontentsline{toc}{section}{Acknowledgements and Declaration of Authorship}
\thispagestyle{plain}

\subsection*{Acknowledgements}

I would like to express my sincere gratitude to my thesis supervisor, Dr. Doan Nhat Quang, for his invaluable guidance, constructive feedback, and continuous support throughout this research. His academic advice and technical insights played an important role in shaping the direction of the study and refining the proposed framework.

I am also grateful to the faculty and staff of the FPT School of Business \& Technology (FSB) for providing a supportive and well-structured academic environment, as well as for the knowledge and experience shared throughout the Master's programme. I would also like to thank my colleagues and peers for their valuable discussions, suggestions, and feedback during the research process.

Finally, I would like to express my heartfelt appreciation to my family for their patience, encouragement, and continued support throughout my graduate studies.

\subsection*{Declaration of Authorship}

I declare that this thesis is my own work and has been prepared in accordance with accepted standards of academic integrity. All sources and materials used in this research have been appropriately acknowledged and cited.

The work presented in this thesis has not been submitted, either in whole or in part, for any other degree or academic qualification at this or any other institution.

I take full responsibility for the content, accuracy, and academic integrity of this thesis.

Nghiem Phan Thanh Tuan -- 2026

\chapter*{Abstract}
\addcontentsline{toc}{chapter}{Abstract}

Vietnamese legal question answering depends on retrieving the specific provisions that support an answer. This thesis studies that retrieval problem under fixed legal corpora, recorded relevance judgments, and article-level evaluation. It asks whether sparse--dense hybrid retrieval improves a lexical BM25 reference, whether adding multilingual-E5-base provides incremental value beyond BM25+BGE-M3, and whether legal-domain adaptation of BGE-M3 transfers from development selection to complete hybrid retrieval.

The study evaluates ALQAC and Zalo under clean and parsed corpus representations, and introduces a Bộ Công an (BCA) question-answer collection derived from published public-service questions and answer citations. The method compares BM25, pretrained BGE-M3, pretrained E5, two-way BM25+BGE-M3 fusion, and three-way BM25+BGE-M3+E5 fusion using NDCG@10 as the primary ranking metric. BCA is evaluated with explicit coverage analysis and a component-aware protocol because many questions share legal targets.

The results show that hybrid retrieval improves the fixed BM25 reference across the reported benchmark conditions. However, the incremental value of E5 is condition-dependent: the clearest positive gain appears on parsed Zalo, while ALQAC shows small negative differences and BCA does not establish a reliable small-component advantage. The BCA analysis also shows that corpus coverage is a substantive limitation: some targets are absent or only partially represented, so ranking improvements alone cannot solve every retrieval failure.

Legal-domain adaptation gives the thesis's most important negative result. Study A: Hybrid-Oriented Adaptation shows substantial development-selection gains, but its corrected held-out hybrid evaluation shows mean NDCG@10 decreases of 0.0520 and 0.0249 for the two-way and three-way hybrids. Study B: Dense-Oriented Adaptation shows that fine-tuning can improve the dense encoder itself, increasing mean NDCG@10 by 0.0189, but this improvement largely disappears after integration into the full hybrids. The implication is that improving a dense retriever does not necessarily improve a complete sparse--dense legal retrieval system.

\textbf{Keywords:} Vietnamese legal retrieval; BM25; BGE-M3; multilingual-E5; hybrid retrieval; contrastive fine-tuning; NDCG; legal-domain adaptation.



\chapter*{List of Abbreviations}
\addcontentsline{toc}{chapter}{List of Abbreviations}

Terminology follows the retrieval literature and the task definitions introduced in Chapters 3--4~\cite{manning2008,chen2024,wang2024,lewis2020}.

\begingroup
\small
\setlength{\tabcolsep}{3pt}
\renewcommand{\arraystretch}{1.18}
\begin{longtable}{>{\raggedright\arraybackslash}p{2.486cm}>{\raggedright\arraybackslash}p{5.073cm}>{\raggedright\arraybackslash}p{7.781cm}}
\toprule
\textbf{Abbreviation} & \textbf{Meaning} & \textbf{Explanation} \\
\midrule
\endfirsthead
\toprule
\textbf{Abbreviation} & \textbf{Meaning} & \textbf{Explanation} \\
\midrule
\endhead
\midrule
\multicolumn{3}{r}{\footnotesize Continued on next page}\\
\endfoot
\bottomrule
\endlastfoot
ALQAC & Automated Legal Question Answering Competition & A benchmark and competition for evaluating legal question answering and information retrieval systems. \\
BCA /\allowbreak{} MPS & Bộ Công an /\allowbreak{} Ministry of Public Security collection & A legal question--document collection derived from public questions and published answers associated with Vietnam's Ministry of Public Security. \\
BM25 & Best Matching 25 lexical ranking function & A term-based retrieval method that ranks documents according to query-term matches, term frequency, and document length. \\
BGE-M3 & BGE-M3 multilingual embedding model & A multilingual embedding model used in this thesis as a dense retrieval component. \\
CV & Cross-validation & An evaluation procedure that partitions data into subsets for model selection and performance estimation. \\
E5 & E5 text-embedding model family & A family of dense embedding models; this thesis evaluates the \texttt{multilin\allowbreak{}gual-e5-\allowbreak{}base} variant. \\
FT /\allowbreak{} PT & Fine-tuned /\allowbreak{} Pretrained & Distinguishes a task-adapted model from its original pretrained version. \\
IR & Information Retrieval & The task of retrieving and ranking documents that are relevant to a user's information need. \\
MAP & Mean Average Precision & A ranking metric that averages precision across relevant results and then across queries. \\
MRR & Mean Reciprocal Rank & A ranking metric based on the reciprocal rank of the first relevant result for each query. \\
NDCG & Normalized Discounted Cumulative Gain & A ranking metric that rewards relevant documents appearing higher in the ranked result list. \\
OOF & Out-of-fold & Predictions or results produced for samples excluded from the corresponding fitting fold. \\
RAG & Retrieval-Augmented Generation & An architecture that retrieves external information before or during text generation. \\
RQ & Research Question & A specific question that the thesis investigates empirically. \\
SD & Standard Deviation & A statistical measure describing dispersion around the mean. \\
\end{longtable}
\endgroup
